\chapter{Discusión y análisis de resultados}
\label{cap:discusion}

Este capítulo interpreta los resultados del Capítulo~\ref{cap:resultados},
los compara con trabajos anteriores sobre el mismo problema y declara sus
limitaciones. Se separa lo que se midió de lo que se interpreta: cuando una
explicación no se comprobó con datos, se dice.

% =====================================================================
\section{Detección}
\label{sec:disc-deteccion}

En un sistema de vigilancia, perder un evento es más grave que dar una falsa
alarma. El evento perdido ya no llega al clasificador, mientras que la falsa
alarma todavía puede terminar en la clase de rechazo. Por eso la meta se puso
en la sensibilidad. Con esa meta, la comparación favorece a PhaseNet: es el
único detector que llega a 0,95 de los volcánicos (119 de 125) y mantiene la
meta cuando se mide en tramos que no se usaron para elegir sus umbrales (0,960
en la rotación de tramos). EQTransformer, con su curva de detección, queda
cerca (0,936) y con menos falsas alarmas (52,7 por hora contra 69,0), pero no
llega, y los intervalos que entrega son demasiado cortos para delimitar los
eventos: solo 31 de 125 quedan bien recortados. STA/LTA es el que da menos
falsas alarmas, pero detecta 0,848 de los volcánicos y su recorte pierde casi
todos los TR.

La ventaja de PhaseNet viene casi entera de la onda S. Con el umbral de la
onda P en 0,40 y el de la S en 0,05, la S detecta 47 volcánicos que la P no
ve, y la P solo 3 que la S no ve. Una explicación posible es que en los LP y
los TR la energía llega de forma gradual, sin el golpe nítido de una onda P de
sismo tectónico, que es con lo que se entrenó la red; no se comprobó. El costo
de depender de la S es doble: PhaseNet necesita un umbral muy bajo, con 69
falsas alarmas por hora, y cuando solo avisa la S el recorte empieza tarde.

La sensibilidad de 0,952 pasa la meta, pero su intervalo de confianza va de
0,899 a 0,978: con 125 eventos no se puede asegurar que el detector supere
0,95 con otro grupo de eventos parecido. Las ventanas de test, que no se
usaron para elegir ningún parámetro, dan una medición independiente: ahí
PhaseNet detecta 845 de 869 eventos principales (0,972), con un intervalo, de
0,959 a 0,981, que queda entero sobre la meta. Las dos mediciones son
coherentes entre sí.

La forma de contar pesa tanto como el detector. Con un aviso a 5\,s del
inicio, PhaseNet detectaba 0,904 de los volcánicos con umbral 0,10, y con
10\,s, 0,968; la diferencia eran solo TR, cuyo aviso cae dentro del evento y
lejos de su inicio. Contar por intersección evita elegir esa tolerancia, que
no viene de ninguna librería, y trata igual a los tres detectores. Su
contrapartida es que cuenta como detectado un evento que el intervalo apenas
toca; por eso la detección se informa siempre junto con el IoU del recorte.

La rotación de tramos mostró que, con más parámetros para elegir, la elección
se ajusta a los eventos con que se hace. En PhaseNet, cambiar también las
ventanas bajaba las falsas alarmas de 69,0 a 48,1 por hora con la traza
completa, pero en tramos nuevos la sensibilidad caía a 0,928. Con 125 eventos
repartidos de forma desigual entre tres tramos la prueba es limitada, pero es
la única forma disponible con estos datos de separar la elección de la
evaluación.

Las falsas alarmas son muchas: unas 69 por hora, casi seis por cada volcánico
detectado. Es un nivel parecido al que informa \textcite{soto2016llaima} para
el detector del volcán Llaima con una estación, cuya FDR era de 0,77 antes de
agregar una etapa que descartaba falsas detecciones y la bajaba a 0,42. En
este trabajo esa tarea queda en manos de la clase de rechazo, que no se
entrenó con falsas alarmas (Sección~\ref{sec:disc-clasificador}).

% =====================================================================
\section{Recorte y espectrogramas}
\label{sec:disc-recorte}

El recorte por energía cumple su papel con VT y LP: en la traza quedan bien
recortados 35 de 36 VT y 49 de 52 LP, y en test, 293 de 300 VT y 281 de 288
LP. Su punto débil es el tremor, que queda bien recortado en 0,65 de los casos
tanto en la traza (24 de 37) como en test (183 de 281). Las mediciones
señalan tres causas. La primera es que el aviso de los TR viene casi siempre
de la onda S y llega en mediana 12,9\,s después del inicio, así que el recorte
empieza tarde. La segunda es que PhaseNet da varios avisos dentro de un mismo
tremor, y el corte en el aviso siguiente lo parte. La tercera es que algunos
TR casi no sobresalen del ruido en FRE y el criterio de energía los termina
antes de tiempo. El corte en el aviso siguiente es un compromiso: suma 12 LP
bien recortados y resta 3 TR.

La duración variable de los eventos se resolvió estirando las columnas del
recorte hasta el ancho de entrada de la red. La ventana fija de 60\,s rinde
casi lo mismo en test (F1 macro de 0,831 a 0,835, contra 0,834 a 0,843 con el
recorte) y en la traza (0,708 a 0,720, contra 0,717 a 0,770). En el promedio,
entonces, el recorte variable no da una ventaja medible. La diferencia
aparece en el tremor cuando el recorte es bueno: con el inicio y el fin del
catálogo, TR se reconoce entre 0,784 y 0,859 con el recorte y entre 0,535 y
0,643 con la ventana fija. Cuando el evento está bien delimitado, conocer su
duración ayuda a reconocer el tremor; el problema de esa clase no está en la
imagen de duración variable, sino en el recorte automático.

Los errores de delimitación se trataron en dos lugares: al etiquetar, aceptando
recortes con IoU de 0,3 o más con su evento, y al entrenar, con la deformación
elástica en el eje del tiempo. Las mediciones por IoU en test muestran que VT
y LP toleran bien esos errores: la sensibilidad de LP casi no cambia con el IoU
y la de VT es incluso algo mayor con IoU bajo que con IoU de 0,7 o más. TR no
los tolera: con IoU entre 0,3 y 0,5 se acierta entre 0,59 y 0,71, y con 0,7 o
más, sobre 0,92.

% =====================================================================
\section{Clasificador}
\label{sec:disc-clasificador}

Las tres redes rinden lo mismo dentro de la variación entre folds: en test, el
F1 macro de VT, LP y TR va de 0,834 a 0,843, con desviaciones de 0,008 a 0,012
y diferencias entre redes menores que la distancia entre el mejor y el peor
modelo de una misma red. Esto le da sentido a la red liviana. EfficientNet-B0,
con unos 5 millones de parámetros en su versión para ImageNet, contra unos 138
millones de VGG16, clasifica más de tres veces más rápido sin perder
desempeño. Con estos resultados no hay razón de desempeño para preferir una
red más grande.

La búsqueda de la configuración, el aumento de datos y la parada temprana
subieron el F1 macro de validación en 0,02 a 0,03. Es una mejora del mismo
orden que la variación entre folds, así que su efecto es modesto. La búsqueda
se hizo con un solo fold y 20 pruebas por red; una búsqueda más amplia podría
encontrar algo mejor, pero no se repitió después de ver test ni la traza,
porque eso sería elegir con los datos de evaluación.

La confusión principal es TR con LP: entre el 13 y el 21\,\% de los TR de test
se clasifican como LP. Las dos clases concentran su energía en frecuencias
bajas, y el tremor se distingue sobre todo por durar más. Las mediciones por
IoU son coherentes con eso: un TR mal recortado, con IoU entre 0,3 y 0,5, se
acierta mucho menos que uno bien recortado.

El punto más débil del clasificador para operar sobre registro continuo son
los recortes sin caja, es decir, los que no tocan ningún evento del catálogo.
Por decisión de método no se usaron para entrenar, porque no se sabe qué son,
y el clasificador manda a la clase de rechazo solo entre un cuarto y un tercio
de ellos. En la traza eso deja entre 60 y 68 recortes por hora llamados
volcánicos sin un evento del catálogo detrás. Una parte podría ser eventos
reales que el catálogo no registró; no se revisó.

% =====================================================================
\section{Sistema completo: dónde se pierde el desempeño}
\label{sec:disc-sistema}

Las mediciones permiten repartir la pérdida del sistema completo. En la traza,
la detección y el recorte dejan bien delimitados 109 de los 125 volcánicos,
lo que pone un techo de 0,872 a la sensibilidad completa; la clasificación la
baja a entre 0,646 y 0,704. La comparación con el recorte de referencia
muestra que lo que cuesta el recorte se concentra en TR: con el inicio y el
fin del catálogo, TR se acierta entre 0,784 y 0,859, y con el sistema, entre
0,432 y 0,519. En VT y LP el recorte automático no cuesta nada medible, y
cuando calza con su evento, clasificar con él rinde igual que con el de
referencia: las diferencias volcánicas van de $-$0,007 a +0,044, con
intervalos que incluyen el cero. Esto responde la pregunta que dio origen al
OE4. Para un clasificador entrenado con recortes automáticos, pasar del recorte
de referencia al automático no degrada la clasificación cuando el recorte
calza; la pérdida viene de los eventos, casi todos TR, que el detector y el
recorte no alcanzan a delimitar.

La otra pérdida, la de VT y la de la clase de rechazo, no viene del recorte.
Con el recorte del catálogo, VT se acierta entre 0,567 y 0,589, igual que con
el automático y muy por debajo de test (0,889 a 0,910). Las comprobaciones
descartaron también que la traza se haya usado para entrenar y que sus
recortes estén peor delimitados que los de test. En la traza, los VT se
confunden con LP y con rechazo, y una cuarta parte del rechazo va a VT. Quedan
dos explicaciones posibles, ninguna medida. Una es que el catálogo de la traza
y el de las ventanas etiquetadas asignen VT, AV e IC con criterios distintos.
La otra es que los tres tramos de la traza correspondan a condiciones del
volcán o del ruido que el entrenamiento no cubre bien, aunque las ventanas de
entrenamiento abarcan de 2014 a 2020 y los tramos son de 2017, 2018 y 2020.
Con 35 VT y 40 eventos de rechazo no se puede separar una de otra; una
revisión de esos eventos por un analista sería la forma directa de hacerlo.

En tiempo, la cadena completa procesa las 9,74 horas de la traza entre 47 y
102\,s según la red, en el procesador de un computador personal y sin tarjeta
gráfica: entre 340 y 740 veces más rápido que el tiempo real. El tiempo no es
una restricción para operar en línea, y la red puede elegirse por desempeño y
no por costo.

% =====================================================================
\section{Comparación con trabajos anteriores}
\label{sec:disc-comparacion}

Los trabajos más cercanos usan el mismo volcán y el mismo conjunto publicado.
\textcite{espinosacurilem2026segmentacion} segmentan con una UNet ventanas de
81,92\,s armadas con la componente vertical de hasta ocho estaciones, y
obtienen en su conjunto de prueba un F1 macro de 0,91 sobre cinco clases (VT
0,92, LP 0,88, TR 0,93, AV 0,90 e IC 0,93) y un IoU medio de 0,88.
\textcite{espinosacurilem2025ood} combinan un VGG16 con un detector de señales
desconocidas y reportan un 84,3\,\% de acierto en las clases volcánicas y un
93,5\,\% en reconocer las no volcánicas. En este trabajo, el F1 macro de VT,
LP y TR en test va de 0,834 a 0,843, con F1 por clase de 0,84 a 0,85 en VT, de
0,86 a 0,87 en LP y de 0,79 a 0,82 en TR.

Las cifras no se pueden comparar directamente, y la diferencia no debe leerse
como una diferencia entre métodos, porque cambian tres condiciones. Aquí se
usa una sola estación y allá hasta ocho, y combinar estaciones mejora el
acierto: \textcite{curilem2016estaciones} midieron una mejora media de 9\,\% en
el volcán Llaima. Aquí cada evento llega recortado por el detector, con su
error de inicio y fin, y allá viene aislado en su ventana. Y los conjuntos de
prueba difieren en composición y en balance de clases. Aun así, la diferencia
mayor está en TR (0,93 contra 0,79 a 0,82), que es justamente la clase que el
recorte automático delimita peor.

El trabajo de \textcite{soto2016llaima}, sobre el volcán Llaima y con una
estación, enfrentó la misma pregunta con un clasificador de máquinas de
vectores de soporte. Entrenado con recortes del detector, rindió 85,05\,\% de
sensibilidad media con esos recortes y 83,69\,\% con los de analistas, y
concluyó que conviene entrenar con los recortes del detector. Aquí el
clasificador también se entrenó con recortes automáticos, y la comparación
pareada en la traza lleva a la misma conclusión con otra técnica y otro
volcán. En operación continua, ese trabajo midió 77,30\,\% de sensibilidad
media de clasificación y 66,04\,\% para el sistema completo, y reporta que,
con una sola estación, ni los analistas superaban el 70\,\% de acierto. Aquí la
sensibilidad completa volcánica va de 0,646 a 0,704. Las medidas no son
iguales, porque allá se multiplica la sensibilidad media de detección por la
de clasificación y aquí se exige que cada evento se detecte, se recorte y se
clasifique bien a la vez; pero en los dos casos el sistema completo queda
claramente por debajo de lo que cada etapa logra por separado.

% =====================================================================
\section{Cumplimiento de los objetivos}
\label{sec:disc-objetivos}

El Cuadro~\ref{tab:disc-objetivos} resume, para cada objetivo específico, lo
que se pidió y lo que se obtuvo. El cumplimiento de cada requerimiento, el
resultado de los criterios de aceptación y el estado final de cada tarea del
backlog están en los Anexos~\ref{sec:anexo-cumplimiento},
\ref{sec:anexo-aceptacion} y~\ref{sec:anexo-estado-backlog}.

{\small
\setlength{\tabcolsep}{4pt}
\begin{longtable}{>{\raggedright\arraybackslash}p{1.0cm}>{\raggedright\arraybackslash}p{4.2cm}>{\raggedright\arraybackslash}p{7.0cm}}
\caption{Cumplimiento de los objetivos específicos.}\label{tab:disc-objetivos}\\
\cabecera{\textbf{Obj.} & \textbf{Qué pide} & \textbf{Qué se obtuvo}}{3}
OE1 & Evaluar STA/LTA, EQTransformer y PhaseNet sobre registro continuo
etiquetado: sensibilidad, falsas alarmas y calidad del recorte (IoU) &
Cumplido. Los tres se evaluaron con la misma forma de contar, con barrido de
sus parámetros y rotación de tramos. PhaseNet es el único que llega a la meta:
0,952 [0,899; 0,978] con 69,0 falsas alarmas por hora; su recorte por energía
deja bien delimitados 0,864 de los volcánicos
(Sección~\ref{sec:res-detectores} y Sección~\ref{sec:res-recorte}). \\
\addlinespace
OE2 & Generar espectrogramas desde los eventos detectados, considerando su
duración variable, los errores de delimitación y el solapamiento &
Cumplido. El recorte se estira a la entrada de la red (duración variable), se
etiqueta con IoU de 0,3 o más y se entrena con deformación elástica (errores
de delimitación), y se corta en el aviso siguiente (solapamiento). TR es la
clase peor recortada: 0,65 en la traza y en test
(Sección~\ref{sec:recorte-espectrogramas} y
Sección~\ref{sec:res-deteccion-test}). \\
\addlinespace
OE3 & Ajustar por transferencia una CNN que clasifique en VT, LP, TR y
rechazo, con protocolo reproducible &
Cumplido con tres redes (VGG16, ResNet50 y EfficientNet-B0), búsqueda de la
configuración y validación cruzada de cinco folds. F1 macro de VT, LP y TR en
test de 0,834 a 0,843. Los datos, las particiones y la semilla quedan fijos;
no se comprobó repitiendo una corrida que dé exactamente lo mismo
(Sección~\ref{sec:res-clasificador}). \\
\addlinespace
OE4 & Evaluar el sistema integrado sobre registro continuo: recorte de
referencia contra automático, por clase y con tiempo de procesamiento &
Cumplido. Sensibilidad completa volcánica de 0,646 a 0,704; frente al recorte
de referencia, el sistema pierde entre 0,056 y 0,098, casi todo en TR, y sin
pérdida cuando el recorte calza. La cadena procesa 9,74 horas en 47 a 102\,s
en CPU (Sección~\ref{sec:res-sistema}). \\
\end{longtable}
}

% =====================================================================
\section{Limitaciones}
\label{sec:disc-limitaciones}

Todo el sistema usa una sola estación, FRE. Los trabajos con varias
estaciones rinden más, y la asociación entre estaciones es la forma habitual de
descartar falsas alarmas, pero queda fuera del alcance de este trabajo.

La traza continua tiene 125 eventos volcánicos, y algunas clases quedan con
pocos recortes (24 TR bien recortados, 35 VT). Los intervalos de confianza
son anchos y varias diferencias no se pueden distinguir del azar. Además, los
umbrales del detector se eligieron con la misma traza con que se evalúa; la
rotación de tramos y las ventanas de test atenúan ese problema, pero no lo
eliminan.

Los detectores se usaron con los pesos publicados, entrenados con sismos
tectónicos, sin ajustarlos con datos volcánicos. El recorte por energía, el
umbral de IoU de 0,3 y los límites del aumento de datos (desplazamiento del
rango de colores de hasta 5\,dB, ancho de 50 a 70\,dB y deformación solo en el
tiempo) son decisiones propias de este trabajo, explicadas en el
Capítulo~\ref{cap:diseno}, y otras elecciones podrían cambiar los resultados.

Los recortes sin caja no se usaron para entrenar, y el clasificador llama
volcánicos a la mayoría de ellos. El catálogo se tomó como la verdad, sin
medir cuánto concuerdan distintos analistas al etiquetar, y esa es una de las
posibles explicaciones de la caída de VT en la traza que quedó sin comprobar.
La búsqueda de la configuración usó un solo fold y 20 pruebas por red.
Finalmente, la evaluación es fuera de línea, sobre registro guardado; no se
probó el sistema recibiendo señal en tiempo real.
